\documentclass[11pt,letterpaper]{article}

\usepackage[margin=1in]{geometry}
\usepackage{amsmath,amssymb}

\linespread{1.08}
\setlength{\parindent}{0pt}
\setlength{\parskip}{0.7em}

\title{%
  Assignment 2: A Bayesian Test of Coin Fairness\\[0.4em]
  {\large CSDS 611 --- Advanced Topics and Methods of Data Science I}%
}
\author{Andrei Cozma}
\date{}

\begin{document}

\maketitle

\section*{Data, hypotheses, and assumptions}

The observed sequence is
\[
  D=(H,H,H,T,T,H),
  \qquad n=6,\qquad k=4\text{ heads and }2\text{ tails}.
\]
Let $\theta$ be the probability of heads on one flip. The hypotheses in the
assignment are
\[
  H_0:\theta=\tfrac12
  \qquad\text{and}\qquad
  H_A:\theta\in[0,1]\text{ is uncertain}.
\]
I assume the flips are independent Bernoulli trials conditional on a shared
$\theta$.

Under $H_A$, I use $\theta\mid H_A\sim\operatorname{Beta}(1,1)$, which is
uniform on $[0,1]$. I also set $P(H_0)=P(H_A)=1/2$ and use the MAP rule, which
chooses the hypothesis with the larger posterior probability.

\section*{Likelihood and model evidence}

For the observed sequence, the likelihood is
\[
  p(D\mid\theta)=\theta^4(1-\theta)^2.
\]
Under $H_0$,
\[
  p(D\mid H_0)
  =\left(\tfrac12\right)^4\left(\tfrac12\right)^2
  =\frac{1}{64}.
\]
Under $H_A$, I average the likelihood over my $\operatorname{Beta}(1,1)$ prior.
Its density is $1$ on $[0,1]$, so
\begin{align*}
  p(D\mid H_A)
    &=\int_0^1 \theta^4(1-\theta)^2\,d\theta \\
    &=\int_0^1 (\theta^4-2\theta^5+\theta^6)\,d\theta \\
    &=\frac15-\frac13+\frac17
     =\frac{1}{105}.
\end{align*}
\section*{Comparison and decision}

The Bayes factor for $H_A$ against $H_0$ is
\[
  \mathrm{BF}_{A0}
  =\frac{p(D\mid H_A)}{p(D\mid H_0)}
  =\frac{1/105}{1/64}
  =\frac{64}{105}\approx0.610.
\]
Thus these data favor $H_0$ over $H_A$ by a factor of $105/64\approx1.64$.
With equal prior probabilities for the hypotheses, Bayes' rule gives
\begin{align*}
  P(H_0\mid D)
    &=\frac{(1/64)(1/2)}{(1/64)(1/2)+(1/105)(1/2)}
     =\frac{105}{169}\approx0.621, \\
  P(H_A\mid D)
    &=\frac{64}{169}\approx0.379.
\end{align*}
The prompt does not specify a prior for $\theta$ under $H_A$, prior
probabilities for the two hypotheses, or a decision rule. Under the choices
stated above, the MAP rule selects $H_0$. The evidence is modest: these six
flips do not establish that the coin is fair.

\end{document}
